% This document was generated by an LLM.

\documentclass{essay}

\newtheorem{conjecture}[theorem]{Conjecture}

\title{The Vocabulary of Mathematical Statements}
\subtitle{Axiom, definition, theorem, lemma, corollary, conjecture---and what the labels actually promise}
\author{}
\date{}

\begin{document}

\maketitle

\section{Why the labels matter}

Open almost any piece of mathematical writing and the text is broken into labeled chunks: an axiom here, a definition there, a theorem, a handful of lemmas leading up to it, a corollary trailing after it. These labels are not decoration. Each one is a small promise about what kind of statement follows, what justification it needs, and how much weight it is meant to carry in the argument.

A definition, for instance, is never something the reader is asked to verify---it is a naming convention, true by stipulation. A theorem, by contrast, is a claim that demands a proof, and the label tells the reader ``this is worth remembering; everything above was building toward it.'' Confusing these categories does not just look sloppy; it hides the actual logical structure of the text; a reader who cannot tell a definition from a claim cannot tell what needs to be checked and what does not.

This essay works through the standard vocabulary, illustrating each term with a single running example: the triangle inequality for real numbers. Along the way it also addresses a point that is easy to miss---some of these labels track genuine logical distinctions (axiom versus theorem), while others track only the author's judgment about importance (theorem versus proposition versus lemma). Knowing which is which is part of learning to read mathematics fluently.

\section{Axioms: statements accepted without proof}

\begin{definition}[Axiom]
An \emph{axiom} (or \emph{postulate}, an older synonym now mostly confined to geometry, as in Euclid's postulates) is a statement accepted as a starting point, without proof, within a given theory. It is not proved because it cannot be proved from anything more basic \emph{inside the theory}---it is part of what defines the theory in the first place.
\end{definition}

An axiom is not a statement that happens to be self-evident; ``self-evident'' is a psychological description, not a mathematical one. Some axioms, like trichotomy below, do feel obvious. Others, like the axiom of choice, are famous precisely because they are not obvious at all, and different mathematicians have historically disagreed about whether to accept them. What makes a statement an axiom is its \emph{role}---it is where the chain of justification stops, by choice, not its intuitive plausibility.

For the running example, take two of the ordering axioms of the real numbers.

\begin{signature}[Trichotomy]
For all real numbers \(a\) and \(b\), exactly one of
\[
a<b,\qquad a=b,\qquad b<a
\]
holds.
\end{signature}

\begin{signature}[Compatibility with addition]
If \(a<b\), then \(a+c<b+c\) for every real number \(c\).
\end{signature}

(These are stated using the \texttt{signature} environment rather than a numbered axiom environment, simply because this essay is not building a full axiomatic theory---but in a text that does build one from scratch, these would appear as \texttt{Axiom 1} and \texttt{Axiom 2}, unproved by design.)

Everything that follows in this essay is a theorem \emph{relative to} axioms like these. That relativity is the point: a theorem is only ever ``proved'' with respect to a stated or assumed set of axioms, and changing the axioms can change what is provable.

\section{Definitions: stipulations, not claims}

\begin{definition}[Definition, informally]
A \emph{definition} introduces a new term or piece of notation as an abbreviation for something already expressible, without asserting anything new about the mathematical world.
\end{definition}

This is the point most often blurred in casual writing. A definition cannot be ``wrong'' in the way a theorem can be wrong---at worst it can be badly chosen (ambiguous, inconsistent with existing usage, or awkward to work with), but it is not a claim that needs a proof. Continuing the running example:

\begin{definition}[Absolute value]
For a real number \(x\), define
\[
|x| =
\begin{cases}
x & \text{if } x \ge 0,\\
-x & \text{if } x < 0.
\end{cases}
\]
\end{definition}

Nothing here is proved, because nothing here is claimed beyond ``this symbol will mean this from now on.'' The mathematical content starts only once we state something \emph{about} \(|x|\), which brings us to the next category.

\section{The proved-statement family}

Theorem, proposition, lemma, corollary, and (informally) claim all denote the same underlying thing: a statement that has been proved from the axioms and definitions in force. They form a single logical category. What distinguishes them from each other is not truth or rigor---all of them, once proved, are equally true---but the author's assessment of \emph{importance} and \emph{role in the argument}. This is a stylistic and expository distinction, not a formal one, and it is worth being explicit about that so the hierarchy is not mistaken for a hierarchy of certainty.

\subsection{Lemma: a stepping stone}

\begin{definition}[Lemma, informally]
A \emph{lemma} is a proved statement whose main purpose is to be used in the proof of something else. Its interest is instrumental rather than intrinsic.
\end{definition}

\begin{lemma}
\label{lem:abs-bounds}
For every real number \(x\),
\[
-|x| \le x \le |x|.
\]
\end{lemma}

\begin{proof}
By trichotomy, either \(x \ge 0\) or \(x<0\).

If \(x\ge 0\), then \(|x|=x\), so \(x\le|x|\) holds with equality. Also \(-|x|=-x\le 0\le x\), using \(x\ge0\).

If \(x<0\), then \(|x|=-x\), so \(x=-|x|\), giving \(-|x|\le x\) with equality. Also \(x<0<-x=|x|\), so \(x\le|x|\).

In both cases \(-|x|\le x\le|x|\).
\end{proof}

Nobody remembers Lemma~\ref{lem:abs-bounds} for its own sake. It exists to make the next proof short.

\subsection{Theorem: the destination}

\begin{definition}[Theorem, informally]
A \emph{theorem} is a proved statement that the author considers a primary result---significant enough to be the point of the section, rather than a means to some other end.
\end{definition}

\begin{theorem}[Triangle inequality]
\label{thm:triangle}
For all real numbers \(a\) and \(b\),
\[
|a+b|\le |a|+|b|.
\]
\end{theorem}

\begin{proof}
By Lemma~\ref{lem:abs-bounds} applied to \(a\) and to \(b\),
\[
-|a|\le a\le|a|
\qquad\text{and}\qquad
-|b|\le b\le|b|.
\]
Adding these two chains of inequalities term by term,
\[
-(|a|+|b|)\le a+b\le |a|+|b|.
\]
The statement \(-M\le y\le M\) is equivalent to \(|y|\le M\) for \(M\ge0\) (this is itself a short lemma, omitted here), so
\[
|a+b|\le |a|+|b|.
\qedhere
\]
\end{proof}

Notice that the lemma did essentially all of the logical work; the theorem's proof is two lines because Lemma~\ref{lem:abs-bounds} had already isolated the useful fact. This is the typical shape of a lemma--theorem pair, and it is why skipping straight to a monolithic theorem statement, with no lemmas, often produces a proof that is correct but hard to follow: the reader cannot see which piece of the argument is doing the work.

\subsection{Proposition: a theorem of secondary rank}

\begin{definition}[Proposition, informally]
A \emph{proposition} is a proved statement of the same logical status as a theorem, but judged by the author to be of lesser importance---interesting in its own right, but not the main event.
\end{definition}

There is no test that distinguishes a proposition from a theorem other than authorial judgment; some textbooks avoid the word ``proposition'' entirely and call everything not clearly a lemma or corollary a theorem. If Theorem~\ref{thm:triangle} had instead appeared in a section devoted to some larger goal---say, proving that \(|\cdot|\) defines a metric on \(\mathbb{R}\)---and the triangle inequality were merely one of three properties checked along the way, many authors would downgrade it to a proposition and reserve ``theorem'' for the metric-space conclusion.

\subsection{Corollary: an easy consequence}

\begin{definition}[Corollary, informally]
A \emph{corollary} is a proved statement that follows quickly---typically in a few lines, with little additional work---from a theorem (or proposition) just established.
\end{definition}

\begin{corollary}[Reverse triangle inequality]
\label{cor:reverse}
For all real numbers \(a\) and \(b\),
\[
\bigl||a|-|b|\bigr| \le |a-b|.
\]
\end{corollary}

\begin{proof}
Write \(a=(a-b)+b\) and apply Theorem~\ref{thm:triangle}:
\[
|a| = |(a-b)+b| \le |a-b|+|b|,
\]
so
\[
|a|-|b|\le |a-b|.
\]
By the same argument with the roles of \(a\) and \(b\) exchanged,
\[
|b|-|a|\le |b-a| = |a-b|.
\]
Combining the two,
\[
\bigl||a|-|b|\bigr| = \max\bigl(|a|-|b|,\ |b|-|a|\bigr) \le |a-b|.
\qedhere
\]
\end{proof}

The label ``corollary'' is doing real communicative work here: it tells the reader, before they read a single line of proof, that this should be easy given what came before, and that if it feels hard they should go back and look for a more direct route through Theorem~\ref{thm:triangle} rather than reinventing the argument from scratch.

\subsection{Claim: a lemma with no life outside its proof}

One more term is common in practice but rarely given a numbered environment: \emph{claim}. A claim is a small proved assertion made \emph{inside} a proof, used once, and never referenced again outside it---for instance, ``Claim: the sequence \((x_n)\) is bounded,'' asserted and dispatched in the middle of a longer argument about convergence. The difference between a claim and a lemma is purely one of scope and permanence: a lemma is promoted out of the proof and given its own label because it will be reused; a claim stays embedded because it will not be. Some authors use ``claim'' and ``lemma'' interchangeably regardless of scope; the distinction described here is a common convention, not a universal rule.

\subsection{Fact: importance suppressed}

Occasionally a true, proved statement is introduced simply as a \emph{fact}, with no numbered environment at all, usually because its proof is either well known, routine, or considered not worth the reader's time to rederive. ``Fact'' carries no information about logical status---a fact is exactly as proved as a theorem---only a signal that the author does not think the proof is worth dwelling on here.

\section{Conjectures: claims without proofs}

\begin{definition}[Conjecture]
A \emph{conjecture} is a mathematical statement proposed as likely true---often on the strength of extensive evidence, partial results, or analogy with proved theorems---but for which no proof is known.
\end{definition}

A conjecture is the mirror image of a theorem: same grammatical form (a precise, falsifiable mathematical statement), opposite epistemic status. A famous example, chosen because it needs no background beyond arithmetic:

\begin{conjecture}[Goldbach, 1742]
Every even integer greater than \(2\) is the sum of two prime numbers.
\end{conjecture}

Goldbach's conjecture has been checked by computer for all even numbers into the quintillions without a single exception, and yet it remains a conjecture, not a theorem, because ``true in every case checked so far'' is not a proof that it is true in every case. The moment someone supplies a valid proof, the statement does not change---only its label does, from conjecture to theorem. This is a useful check on the vocabulary: the words describe our \emph{knowledge} about a statement, not the statement itself, which was either always true or always false regardless of when, or whether, anyone proves it.

(A conjecture that is later proved false is usually just called a disproved conjecture, or is remembered by a counterexample; a well-known instance is Euler's sum of powers conjecture, refuted by a single counterexample found in 1966 after standing since 1769.)

\section{Supporting apparatus}

A handful of further terms round out the vocabulary, describing not statements but the material that surrounds them.

\begin{itemize}
\item \textbf{Proof.} The argument establishing that a theorem (or lemma, proposition, corollary) is a logical consequence of the axioms, definitions, and previously proved statements. Everything above this section contained proofs; a theorem without one is, until supplied, only a conjecture wearing a theorem's label.
\item \textbf{Example.} A specific instance illustrating a definition or a theorem, used for concreteness rather than to establish anything general.
\item \textbf{Exercise.} A statement left for the reader to prove, typically routine enough that spelling out the proof would not teach anything the reader cannot work out themselves.
\item \textbf{Remark.} A comment that adds context, intuition, or a connection to other ideas, but asserts nothing that requires proof. The observation in the previous section, that a proved conjecture changes label but not truth value, is exactly this kind of remark.
\item \textbf{Note.} Used interchangeably with remark in most texts (including this one's document class), sometimes with a narrower connotation of a brief clarification rather than a discursive aside.
\end{itemize}

\begin{remark}
Two further terms are worth knowing by sight even though they are rare today. A \emph{porism}, in Euclid's usage, was a statement falling somewhere between a theorem and a corollary---a byproduct of a construction, noted because it seemed interesting rather than sought directly; the word has essentially fallen out of use. A \emph{scholium} was a discursive comment appended to a proposition, playing much the same role ``remark'' plays now. Neither term appears in modern usage outside of historical texts and the occasional deliberately old-fashioned title.
\end{remark}

\section{A word that names a technique, not a statement}

One more term deserves a mention because it is easy to place in the wrong category: \emph{principle}, as in the ``pigeonhole principle'' or the ``principle of mathematical induction.'' A principle is not a separate rank alongside axiom and theorem; it is simply a name attached to a statement (or a proof technique built on one) that has become so widely used that it acquired an informal title of its own. The induction principle, for instance, is an axiom in the Peano axiomatization of the natural numbers, but a \emph{theorem} in set-theoretic constructions of \(\mathbb{N}\) where it is derived from well-ordering. Whether a ``principle'' is an axiom or a theorem depends entirely on which foundational framework it is being stated in---the word ``principle'' itself carries no information about that either way.

\section{The hierarchy is judgment, the categories are not}

It is worth separating two very different distinctions this essay has drawn, because collapsing them is the most common source of confusion.

The distinction between an \textbf{axiom} and everything else is a genuine logical one: axioms are unproved by design, and every theorem, lemma, proposition, and corollary is proved from them. Likewise, the distinction between a \textbf{definition} and a proved statement is genuine: a definition asserts nothing and needs no proof; a theorem asserts something and needs one. And the distinction between a \textbf{conjecture} and a theorem is genuine and, unlike the others, can actually change over time as proofs are found.

By contrast, the distinction \emph{among} theorem, proposition, lemma, and corollary is not logical at all---it is the author's running commentary on the shape of their own argument: this is the destination (theorem), this is worth having but secondary (proposition), this exists to serve something else (lemma), this falls out for free (corollary). A statement does not have an intrinsic rank in this second hierarchy; the same fact, moved into a different paper with a different goal, could easily change labels. Theorem~\ref{thm:triangle} above is a theorem in this essay because the triangle inequality was the point of writing it; in a real analysis textbook building up the theory of metric spaces, it typically appears as one proposition among several, on the way to a theorem about \(\mathbb{R}\) being a metric space.

Reading mathematics well means tracking both hierarchies at once: trusting the first one absolutely (an axiom really is unproved; a theorem really is proved from it), and reading the second one as a guide to emphasis rather than a claim about certainty.

\end{document}
